% Replace authoring prompts with a case-specific scientific argument.
% Visible filenames, paths, scripts, and routine names are prohibited.
% Explain terms at first use in prose or beside equations; no glossary section.


\documentclass[../report/case_report.tex]{subfiles}
\begin{document}

\section{Verification purpose}
\textbf{Objective:} State the physical or numerical behavior being verified,
its importance to ELMFIRE, controlled processes, and the precise comparison
question. Distinguish verification from validation against observations.


\section{Mathematical formulation and numerical implementation}
List assumptions concisely:
\begin{itemize}[nosep]
\item Homogeneous fuel; no spotting; steady uniform wind.
\item 2D structured grid, constant \(\Delta x = \Delta y\).
\item Boundary conditions: specify.
\end{itemize}


\section{Derivation of expected behavior}
Derive the reference response for the prescribed conditions. Identify whether
it is analytical, independently calculated, or based on conservation, symmetry,
convergence, or a limiting response. Restate and cite necessary external theory.

\section{Verification metrics and acceptance criteria}
Define each metric, its selection rationale, region, mask, normalization, time
and ensemble selection, and justified tolerance. State the complete decision
rule. Missing evidence is NOT EVALUABLE, never evidence of PASS.

\section{Simulation configuration and predicted outcome}
Describe physical and numerical choices by scientific meaning and units, not
configuration identifiers. State predicted responses before measured results.


\subsection{Input data}
Describe each input's physical quantity, units, value or spatial distribution,
domain extent, grid dimensions and resolution, coordinate reference system,
and time basis. Explain fuels, moisture, wind height and direction, terrain,
canopy or structures, ignition geometry and timing, buffers, missing and
nonburnable cells, transformations, and random distributions/seeds as applicable.
Cite external data sources. For unknown details state ``not documented in the
available case materials.'' Never invent a reproducibility claim.


\subsection{Numerical controls}
Describe mesh resolution, timestep, the Courant--Friedrichs--Lewy (CFL)
stability condition, and the level-set discretization.


\section{Actual simulation results}
Present available evidence, completeness, and the ELMFIRE release or source
revision if known. Describe physical output quantities and their unique
time/member selection. Distinguish simulation from prediction and observation.


\section{Calculated metrics and verification decision}
Present expected limits, calculated values, and component decisions from the
established metric calculations. Preserve PASS, FAIL, CHARACTERIZED, and
NOT EVALUABLE. Report preparation cannot establish scientific agreement.


\section{Technical interpretation}
Explain agreement or disagreement using the derivation, resolution, assumptions,
uncertainty, and available evidence. Do not merely restate the results table.


Captions identify the field, units, conditions, spatial and temporal context,
and scientific purpose. Define abbreviations and keep labels readable. Do not
display filenames, operational paths, or software-development terminology.


\section{Scope and limitations}
State what the evidence demonstrates and what remains unverified. Explain
limits on extrapolation beyond the prescribed conditions.


\end{document}
